\section{Results}
\label{sec:results}

\subsection{Main results}
\label{sec:results:main}

Table~\ref{tab:main} presents the six-method comparison across four
protocols. Three observations drive the rest of this section.

% ---- Main results table ----
\begin{table*}[t]
\centering
\small
\begin{tabular}{lcccc}
\toprule
\textbf{Method} & \textbf{GoEmotions LODO} & \textbf{ISEAR LODO} & \textbf{WASSA-21 LODO} & \textbf{Mixed} \\
\midrule
CE source-only         & $0.300 \pm 0.007$ & $0.517 \pm 0.014$ & $0.464 \pm 0.020$ & $0.710 \pm 0.004$ \\
DANN ($\lambda_{\max}{=}0.5$) & $0.291$\textsuperscript{$\dagger$} & $0.485$\textsuperscript{$\dagger$} & $0.473$\textsuperscript{$\dagger$} & $\approx 0.71$\textsuperscript{$\dagger$} \\
CDAN                   & $0.294 \pm 0.015$ & $0.512 \pm 0.000$ & $0.498 \pm 0.014$ & $0.714 \pm 0.005$ \\
\textbf{CE + Focal}    & $0.296 \pm 0.002$ & $0.508 \pm 0.018$ & $\mathbf{0.516 \pm 0.005}$\textsuperscript{$\ast$} & $0.690 \pm 0.011$\textsuperscript{$\ast$} \\
DANN + Focal           & $0.308$\textsuperscript{$\dagger$} & $0.498$\textsuperscript{$\dagger$} & $0.501 \pm 0.016$ & $0.669$\textsuperscript{$\dagger$} \\
CDAN + Focal           & $0.290$\textsuperscript{$\dagger$} & $0.481$\textsuperscript{$\dagger$} & $0.483 \pm 0.013$ & $0.677$\textsuperscript{$\dagger$} \\
\bottomrule
\end{tabular}
\caption{Test macro-F1 (mean $\pm$ standard deviation across three
seeds, DeBERTa-v3-base). $^\ast$ indicates $p < 0.05$ vs the
CE-source-only baseline in the same protocol under a 1000-resample
paired bootstrap. $^\dagger$ indicates single-seed (seed 42) result;
these cells are not included in significance tests. The largest
positive deviation from CE is \textbf{CE + Focal on WASSA-21 LODO}
(\focalwassagain{} F1 points, $p = \focalwassap$).}
\label{tab:main}
\end{table*}

\paragraph{Finding 1: focal alone wins.}
CE + Focal achieves \focalwassagain{} F1 points over CE source-only on
WASSA-21 LODO ($0.516 \pm 0.005$ vs $0.464 \pm 0.020$, paired bootstrap
$p = \focalwassap$). This is the largest cross-dataset improvement we
observe under any of the six methods, on any of the three LODO
targets. The result is counter-intuitive: of the six methods, the
simplest non-adversarial variant produces the strongest gain.

\paragraph{Finding 2: adversarial layers do not amplify focal.}
On the same target (WASSA-21), DANN + Focal achieves $0.501 \pm 0.016$
(+1.8\,$\sigma$ vs CE, bootstrap $p \approx 0.11$, marginal) and
CDAN + Focal achieves $0.483 \pm 0.013$ ($p \approx 0.13$, n.s.).
\emph{Adding} an adversarial component to focal training does not
amplify the gain; in both cases it reduces it. We discuss why in
Section~\ref{sec:discussion:combinations}.

\paragraph{Finding 3: two LODO targets are method-invariant.}
On GoEmotions LODO, all six methods produce test macro-F1 between
$0.290$ and $0.308$ — within $1.5$ standard deviations of CE's
seed-to-seed variance. On ISEAR LODO, all six methods fall between
$0.481$ and $0.517$, again within seed variance. No pair of methods on
these two targets reaches $p < 0.1$ under paired bootstrap. This
contrasts sharply with the WASSA-21 panel, where four of six methods
exceed CE by $1.5$+ standard deviations.

\subsection{Pairwise bootstrap significance}
\label{sec:results:bootstrap}

% ---- Bootstrap p-values table (subset of pairwise_pvalues.csv) ----
\begin{table}[t]
\centering
\small
\begin{tabular}{llc}
\toprule
\textbf{Comparison} & \textbf{Target} & \textbf{p-value} \\
\midrule
\multicolumn{3}{l}{\emph{Significant positive (focal vs CE):}} \\
CE + Focal vs CE source-only & WASSA LODO & \textbf{0.014} \\
\midrule
\multicolumn{3}{l}{\emph{Significant negative (focal harms Mixed):}} \\
CE + Focal vs CE source-only & Mixed & \textbf{$<$0.001} \\
CE + Focal vs CDAN          & Mixed & \textbf{$<$0.001} \\
\midrule
\multicolumn{3}{l}{\emph{Marginal (within $0.05 < p < 0.15$):}} \\
CDAN vs CE source-only      & WASSA LODO & 0.076 \\
DANN + Focal vs CE source-only & WASSA LODO & 0.106 \\
\midrule
\multicolumn{3}{l}{\emph{Method-invariant targets:}} \\
CE + Focal vs CE source-only & GoEmotions LODO & 0.388 \\
CE + Focal vs CE source-only & ISEAR LODO & 0.436 \\
CDAN vs CE source-only      & GoEmotions LODO & 0.14 \\
CDAN vs CE source-only      & ISEAR LODO & 0.72 \\
\midrule
\multicolumn{3}{l}{\emph{Mixed protocol stability:}} \\
CDAN vs CE                   & Mixed & 0.246 \\
\bottomrule
\end{tabular}
\caption{Selected paired bootstrap $p$-values (1000 resamples, seed
42). The complete pairwise table is in
Appendix~\ref{app:bootstrap}.}
\label{tab:bootstrap}
\end{table}

\subsection{Per-class analysis on WASSA-21 LODO}
\label{sec:results:perclass}

Row-normalized confusion matrices on the WASSA-21 test set (seed 42)
reveal the mechanism of focal loss's gain
(Figure~\ref{fig:confusion}). CE source-only achieves zero recall on
\emph{disgust} (the rarest class in the source training distribution)
and only $0.30$ recall on \emph{fear}. CE + Focal recovers these to
$0.12$ and $0.45$ respectively, at the cost of a small reduction on
\emph{surprise} ($0.65 \rightarrow 0.59$). The +5 F1-point macro
improvement is driven primarily by these two rare-class recoveries.

Figure~\ref{fig:perclassf1} makes the same mechanism explicit in
terms of per-class F1: the largest bar-height gains for CE + Focal
over CE source-only are on the rare classes \emph{disgust} and
\emph{fear}, while the high-frequency classes (\emph{joy},
\emph{sadness}) are essentially unchanged. This confirms that focal
loss's contribution is rare-class recovery rather than a uniform
improvement across the label space.

\begin{figure}[t]
\centering
\includegraphics[width=\linewidth]{../outputs/figures/perclass_f1_wassa.pdf}
\caption{Per-class test F1 on the WASSA-21 LODO test set (seed 42).
CE + Focal's gains are concentrated on the rare classes (disgust,
fear); per-class deltas are annotated below each class.}
\label{fig:perclassf1}
\end{figure}

\begin{figure*}[t]
\centering
\includegraphics[width=0.95\linewidth]{../outputs/figures/confusion_matrix_wassa.pdf}
\caption{Row-normalized confusion matrices on the WASSA-21 LODO test
set (seed 42). Left: CE source-only. Right: CE + Focal. Focal loss
recovers rare classes (disgust, fear) at a small cost on surprise.}
\label{fig:confusion}
\end{figure*}

\subsection{Validation-to-test gap pattern}
\label{sec:results:gap}

Figure~\ref{fig:gap} shows validation and test macro-F1 side-by-side
for each method on each LODO target. The
gap is method-invariant on GoEmotions ($0.42$--$0.47$) and ISEAR
($0.18$--$0.22$): every method overfits its source validation set by
the same amount. On WASSA-21, CE + Focal shrinks the gap from $0.26$
(CE source-only) to $0.18$, the same protocol where it improves test
F1 the most. This pattern is consistent with focal loss reducing
source-domain overfitting under class imbalance, rather than
improving raw classification capacity.

\begin{figure*}[t]
\centering
\includegraphics[width=0.95\linewidth]{../outputs/figures/val_test_gap.pdf}
\caption{Validation vs test macro-F1 by method and LODO target.
Gap is method-invariant on GoEmotions and ISEAR; CE + Focal closes
the gap on WASSA-21.}
\label{fig:gap}
\end{figure*}

\subsection{DeBERTa-v3-large robustness validation}
\label{sec:results:large}

% ---- Large model 5-seed table ----
\begin{table}[t]
\centering
\small
\begin{tabular}{lccc}
\toprule
\textbf{Configuration} & \textbf{n} & \textbf{mean} & \textbf{std} \\
\midrule
CE source-only (base)  & 3 & 0.464 & 0.020 \\
CE + Focal (base)      & 3 & \textbf{0.516} & 0.005 \\
CE + Focal (large)     & 5 & 0.492 & 0.029 \\
\bottomrule
\end{tabular}
\caption{Robustness check on the single strongest cell (CE + Focal,
WASSA-21 LODO). DeBERTa-v3-large 5-seed mean is intermediate between
base focal and base CE; Welch's $t$-test fails to reject either
hypothesis ($p \approx 0.13$ vs base focal, $p \approx 0.16$ vs base
CE). We interpret this as inflated seed variance under the small
WASSA-21 test set rather than as evidence against focal at scale; see
Section~\ref{sec:discussion:scale}.}
\label{tab:large}
\end{table}

\paragraph{Variance, not mean, is the qualitative change.}
Source-validation macro-F1 across the five large-model seeds remains
tightly clustered (range $[0.677, 0.683]$, std $0.002$). The
variance is concentrated entirely in held-out target macro-F1, which
ranges from $0.462$ to $0.534$. This pattern — stable source learning,
unstable target generalization — indicates that the larger model's
extra capacity is being absorbed by source-domain overfitting at
some random initializations and not at others, rather than the
focal-loss mechanism failing systematically.

\paragraph{What this means for the headline claim.}
The base-model focal-loss gain ($p = \focalwassap$ under paired
bootstrap) is the primary contribution of this paper. The
DeBERTa-large 5-seed result is reported as a robustness check; it
neither confirms nor disconfirms the gain at the larger scale, and we
report this transparently rather than over-claiming. Reporting both
the gain and its scale-dependent variance increase is, we argue, more
honest than running additional seeds until a positive result is
recovered, given the small WASSA-21 test set ($n = 164$) sets a hard
ceiling on bootstrap resolution at large scale.
